% Article VII. Incorporation of pre-existing results, 10-09-2026.
% Complete source owners preserved in fuentes_conservadas:
% 17_pantalla_cubica_soldadura.tex (P), 18_clausura_energetica.tex (E).
% P: incidence, quotient, cellular soldering, and refinement.
% E: response, uniqueness, minimal extension, and tracial expectation.
% Previous dependencies of VII: APP--TRIT--TPK kernel; cubic realization
% of the route, transport cocycle, and compression q_vis=5/9.
% This fragment alone does not constitute the self-contained article.
% The screen identity is distinguished from the material spin contraction.

\subsection{Cubic incidence and boundary soldering}
\label{vii:sec:pantalla-respuesta}

The cubic realization of transport preserves six oriented incidences.
Its geometric quotient, Lorentzian form, and energetic response
are obtained from the same incidence matrix. Transport originates in
the enriched APP--TRIT--TPK route: the face, sign, carry, and
memory accompany each edge. This section develops the
linear operation acting on that boundary representation.

\subsubsection{Cubic incidence and the rank-four quotient}
\label{vii:sec:cociente-pantalla}

Let \(H_F=\mathbb R^F\), with
\(F=\{(i,\epsilon):1\leq i\leq3,\ \epsilon\in\{-1,1\}\}\),
its Euclidean inner product, and its basis \(e_{i,\epsilon}\).
Face opposition is the orthogonal involution
\(Re_{i,\epsilon}=e_{i,-\epsilon}\). Define
\[
 u=\sum_{i,\epsilon}e_{i,\epsilon},\qquad
 P_u=\frac{uu^{\mathsf T}}6,\qquad P_-=\frac{I-R}{2},\qquad
 P_0=\frac{I+R}{2}-P_u,\qquad P_\square=P_u+P_-.
\]
The vertices are \(V=\{-1,1\}^3\). The incidence matrix
\(M:\mathbb R^V\longrightarrow H_F\) is determined by
\[
 M_{(i,\epsilon),\nu}=\mathbf1_{\{\nu_i=\epsilon\}}.
\]

\begin{proposition}[Incidence decomposition]
\label{vii:prop:incidencia}
One has
\[
 MM^{\mathsf T}=12P_u+4P_-,
 \qquad
 \ker M^{\mathsf T}=P_0H_F.
\]
Consequently, the quotient
\(Q_\square=Q=H_F/P_0H_F\), orthogonally identified with
\(P_\square H_F\), has dimension four.
\end{proposition}

\begin{proof}
A face contains four vertices. Opposite faces have no common vertices,
and two faces on different axes have two common vertices.
Writing \(\mathbf J=uu^{\mathsf T}\) gives
\[
 MM^{\mathsf T}=4I+2(\mathbf J-I-R)
 =2(I-R)+2\mathbf J=4P_-+12P_u.
\]
The three projectors \(P_u,P_-,P_0\) are orthogonal and sum to \(I\);
their ranks are one, three, and two, respectively. The identity
\(\|M^{\mathsf T}x\|^2=\langle x,MM^{\mathsf T}x\rangle\)
identifies the stated kernel. The complementary rank is four.
\end{proof}

On \(Q\), the form
\[
 B=P_--P_u
\]
has signature \((3,1)\). Time orientation is fixed by \(u\).
With
\[
 t=\frac{u}{\sqrt6},\qquad
 d_i=\frac{e_{i,+}-e_{i,-}}{\sqrt2},\qquad
 W=\sqrt6P_u+\sqrt2P_-,
\]
one has \(We_{i,\epsilon}=t+\epsilon d_i\). These six vectors
are null for \(B\), whereas
\[
 3t+\nu_1d_1+\nu_2d_2+\nu_3d_3
\]
has squared norm \(-6\). The same incidence satisfies
\(MM^{\mathsf T}=2W^*W\) on \(Q\). The proper rotations of the cube
commute with the projectors, preserve both forms, and fix \(t\).
These identities determine the metric of the cubic representation.

\subsubsection{Cellular soldering and memory transport}
\label{vii:sec:soldadura-celular}

For each oriented edge \(a\), the reading of the route provides
the face \(\lambda(\underline a)\), the sign
\(\sigma(a)\), the source and target fibers, and the transport
\(U(a):Q_{s(a)}\longrightarrow Q_{t(a)}\).
Reversal satisfies \(U(\bar a)=U(a)^{-1}\).
Boundary half-edges preserve their formal reversal and their incidence
on the boundary. With \(\ell_m=\ell_0\,9^{-m}>0\), soldering is
\begin{equation}
 \beta_m(a)=\ell_m\sigma_m(a)
              n_{\lambda(\underline a)}^{s(a)},\qquad
 n_{i,\epsilon}=t+\epsilon d_i.
 \label{vii:eq:soldadura-arista}
\end{equation}
The face and sign are transported together with the route memory.
The reversal convention requires
\begin{equation}
 \beta_m(\bar a)=-U_m(a)\beta_m(a).
 \label{vii:eq:inversion-soldadura}
\end{equation}
Applying this relation twice recovers \(\beta_m(a)\);
the carry record separately preserves the ordered composition.

\begin{proposition}[Naturality of soldering under refinement]
\label{vii:prop:refinamiento-soldadura}
Let \(a=a_1\cdots a_9\) be a compatible refinement of an edge.
Suppose that the nine successor incidences, transported to the predecessor fiber,
preserve the face and sign, and that
\(\ell_{m+1}=\ell_m/9\). Let
\(\mathcal U_j:Q_{m,s(a)}\longrightarrow Q_{m+1,s(a_j)}\)
be the accumulated transport, including the refinement identification
and the path to the source of successor \(j\). Then
\[
 \beta_m(a)=
 \sum_{j=1}^{9}\mathcal U_j^{-1}\beta_{m+1}(a_j).
\]
The equality is compatible with reversal and with concatenation
of the memory record.
\end{proposition}

\begin{proof}
Each summand, expressed in the source fiber of the predecessor, is
\((\ell_m/9)\sigma_m(a)n_{\lambda(\underline a)}\).
The sum of nine terms reproduces the predecessor soldering.
Reversal reverses the order of the transports and replaces each
by its inverse; relation \eqref{vii:eq:inversion-soldadura} supplies
the overall sign. Memory is concatenated in the same order, so
refinement preserves both the vector representation and its
enriched record.
\end{proof}
